\chapter{2008年浙江卷（文）}

\section{单选题}

\begin{problem}
已知集合 \(A = \{x \mid x > 0\}\)，\(B = \{x \mid -1 \leqslant x \leqslant 2\}\)，则 \(A \cup B =\)（\quad）

\choices
    {\( \{x \mid x \geqslant -1\} \)}
    {\(\{x\mid x\leqslant 2\}\)}
    {\( \{x \mid 0 < x \leqslant 2\} \)}
    {\(\{x\mid -1\leqslant x\leqslant 2\}\)}
\end{problem}

\begin{answer}
A
\end{answer}

\begin{solution}
由 \(B=\{x\mid -1\leqslant x\leqslant2\}\)，可知 \(B\) 中的每个数都满足 \(x\geqslant-1\)，而 \(A=\{x\mid x>0\}\) 中的数也都满足 \(x\geqslant-1\)，所以 \(A\cup B\subseteq\{x\mid x\geqslant-1\}\). 反之，当 \(-1\leqslant x\leqslant2\) 时 \(x\in B\)，当 \(x>2\) 时 \(x\in A\)，故每个满足 \(x\geqslant-1\) 的数都属于 \(A\cup B\). 因而 \(A\cup B=\{x\mid x\geqslant-1\}\)，选择 A.
\end{solution}

\begin{problem}
函数 \( y = (\sin x + \cos x)^2 + 1 \) 的最小正周期是（\quad）

\choices
    {\(\frac{\pi}{2}\)}
    {\(\pi\)}
    {\(\frac{3\pi}{2}\)}
    {\(2\pi\)}
\end{problem}

\begin{answer}
B
\end{answer}

\begin{solution}
展开并使用 \(\sin^2x+\cos^2x=1\) 及 \(2\sin x\cos x=\sin2x\)，得 \(y=1+2\sin x\cos x+1=2+\sin2x\). 因为 \(\sin2x\) 的周期为 \(\pi\)，且不存在比 \(\pi\) 更小的正数 \(T\) 使 \(\sin(2x+2T)=\sin2x\) 对所有 \(x\) 成立，所以原函数的最小正周期为 \(\pi\)，选择 B.
\end{solution}

\begin{problem}
已知 \(a\)，\(b\) 都是实数. 那么“\(a^2 > b^2\)”是“\(a > b\)”的（\quad）

\choices
    {充分而不必要条件}
    {必要而不充分条件}
    {充分必要条件}
    {既不充分也不必要条件}
\end{problem}

\begin{answer}
D
\end{answer}

\begin{solution}
先验证 \(a^2>b^2\) 不能推出 \(a>b\)：取 \(a=-2\)，\(b=1\)，则 \(a^2=4>b^2=1\)，但 \(a<b\). 再验证 \(a>b\) 不能推出 \(a^2>b^2\)：取 \(a=1\)，\(b=-2\)，则 \(a>b\)，但 \(a^2=1<b^2=4\). 两个条件都不能推出另一个条件，所以“\(a^2>b^2\)”是“\(a>b\)”的既不充分也不必要条件，选择 D.
\end{solution}

\begin{problem}
已知 \(\{a_{n}\}\) 是等比数列，\(a_{2} = 2\)，\(a_{5} = \frac{1}{4}\)，则公比 \(q =\)（\quad）

\choices
    {\(-\frac{1}{2}\)}
    {\(-2\)}
    {\(2\)}
    {\(\frac{1}{2}\)}
\end{problem}

\begin{answer}
D
\end{answer}

\begin{solution}
等比数列满足 \(a_5=a_2q^{5-2}\)，所以 \(2q^3=\frac14\)，即 \(q^3=\frac18\). 实数立方函数在 \(\R\) 上单调递增，故 \(q=\frac12\)，选择 D.
\end{solution}

\begin{problem}
已知 \(a \geqslant 0\)，\(b \geqslant 0\)，且 \(a + b = 2\)，则（\quad）

\choices
    {\(ab \leqslant \frac{1}{2}\)}
    {\(ab \geqslant \frac{1}{2}\)}
    {\(a^2 + b^2 \geqslant 2\)}
    {\(a^2 + b^2 \leqslant 3\)}
\end{problem}

\begin{answer}
C
\end{answer}

\begin{solution}
由 \((a-b)^2\geqslant0\)，得 \((a+b)^2\geqslant4ab\)，从而 \(4\geqslant4ab\)，即 \(ab\leqslant1\). 又由 \(a^2+b^2=(a+b)^2-2ab=4-2ab\)，可得 \(a^2+b^2\geqslant2\)，所以选项 C 恒成立. 取 \(a=b=1\) 时 \(ab=1>\frac12\)，故 A 不恒成立；取 \(a=0\)，\(b=2\) 时 \(ab=0<\frac12\)，故 B 不恒成立，同时 \(a^2+b^2=4>3\)，故 D 不恒成立. 因此选择 C.
\end{solution}

\begin{problem}
在 \((x - 1)(x - 2)(x - 3)(x - 4)(x - 5)\) 的展开式中，含 \(x^4\) 的项的系数是（\quad）

\choices
    {\(-15\)}
    {\(85\)}
    {\(-120\)}
    {\(274\)}
\end{problem}

\begin{answer}
A
\end{answer}

\begin{solution}
每个因式 \((x-k)\) 只含一次 \(x\). 要得到次数为 \(4\) 的项，必须从五个因式中有四个取 \(x\)，另一个取常数项 \(-k\). 因而 \(x^4\) 的系数为 \((-1)+(-2)+(-3)+(-4)+(-5)=-15\)，选择 A.
\end{solution}

\begin{problem}
在同一平面直角坐标系中，函数 \( y = \cos \left( \frac{x}{2} + \frac{3\pi}{2} \right) \)（\(x \in [0, 2\pi]\)）的图象和直线 \( y = \frac{1}{2} \) 的交点个数是（\quad）

\choices
    {\(0\)}
    {\(1\)}
    {\(2\)}
    {\(4\)}
\end{problem}

\begin{answer}
C
\end{answer}

\begin{solution}
由 \(\cos\left(t+\frac{3\pi}{2}\right)=\sin t\)，令 \(t=\frac{x}{2}\)，原函数为 \(y=\sin\frac{x}{2}\). 当 \(x\in[0,2\pi]\) 时，\(\frac{x}{2}\in[0,\pi]\)，所以交点的横坐标满足 \(\sin\frac{x}{2}=\frac12\). 在 \([0,\pi]\) 内，方程的两个解为 \(\frac{x}{2}=\frac\pi6\) 和 \(\frac{x}{2}=\frac{5\pi}{6}\)，对应 \(x=\frac\pi3\) 和 \(x=\frac{5\pi}{3}\). 因而交点有 \(2\) 个，选择 C.
\end{solution}

\begin{problem}
若双曲线 \(\frac{x^2}{a^2} - \frac{y^2}{b^2} = 1\) 的两个焦点到一条准线的距离之比为 \(3:2\)，则双曲线的离心率是（\quad）

\choices
    {\(3\)}
    {\(5\)}
    {\(\sqrt{3}\)}
    {\(\sqrt{5}\)}
\end{problem}

\begin{answer}
D
\end{answer}

\begin{solution}
设双曲线的半实轴长、焦半距和离心率分别为 \(a\)，\(c\)，\(e\)，则 \(c=ae\)，且双曲线的右准线为 \(x=\frac{a}{e}\). 因为双曲线的离心率满足 \(e>1\)，所以 \(c>\frac{a}{e}\). 两个焦点 \((c,0)\)，\((-c,0)\) 到右准线的距离分别为 \(c-\frac{a}{e}\) 和 \(c+\frac{a}{e}\). 令 \(t=\frac{a/e}{c}=\frac1{e^2}\)，则 \(0<t<1\)，并且由题意
\[\frac{1+t}{1-t}=\frac32\]
解得 \(2+2t=3-3t\)，所以 \(t=\frac15\)，即 \(e^2=5\). 由于离心率为正，故 \(e=\sqrt5\)，选择 D.
\end{solution}

\begin{problem}
对两条不相交的空间直线 \(a\) 与 \(b\)，必存在平面 \(\alpha\)，使得（\quad）

\choices
    {\(a \subset \alpha, b \subset \alpha\)}
    {\(a \subset \alpha\)，\(b \parallel \alpha\)}
    {\(a \perp \alpha\)，\(b \perp \alpha\)}
    {\(a \subset \alpha\)，\(b \perp \alpha\)}
\end{problem}

\begin{answer}
B
\end{answer}

\begin{solution}
若 \(a\)，\(b\) 异面，过直线 \(a\) 上任一点作直线 \(b'\parallel b\)，则相交直线 \(a\) 与 \(b'\) 确定平面 \(\alpha\). 因为 \(b'\subset\alpha\)，所以 \(b\) 的方向平行于平面 \(\alpha\). 若 \(b\subset\alpha\)，则 \(a\)，\(b\) 共面且不相交，从而应有 \(a\parallel b\)，与异面矛盾，所以 \(b\not\subset\alpha\)，即 \(b\parallel\alpha\).

若 \(a\parallel b\)，在所有经过 \(a\) 的平面中取一个不包含 \(b\) 的平面 \(\alpha\). 由于 \(b\) 的方向与 \(a\) 平行，故其方向平行于 \(\alpha\)；又 \(b\not\subset\alpha\) 且 \(a\)，\(b\) 不相交，所以 \(b\) 与 \(\alpha\) 无公共点，从而 \(b\parallel\alpha\). 因此总能使 \(a\subset\alpha\)，\(b\parallel\alpha\)，选择 B.
\end{solution}

\begin{problem}
若 \(a \geqslant 0\)，\(b \geqslant 0\)，且当 \(\begin{cases}
x \geqslant 0,\\
y \geqslant 0,\\
x + y \leqslant 1
\end{cases}\) 时，恒有 \(ax + by \leqslant 1\)，则以 \(a\)，\(b\) 为坐标的点 \(P(a, b)\) 所形成的平面区域的面积是（\quad）

\choices
    {\( \frac{1}{2} \)}
    {\( \frac{\pi}{4} \)}
    {\(1\)}
    {\(\frac{\pi}{2}\)}
\end{problem}

\begin{answer}
C
\end{answer}

\begin{solution}
变量 \((x,y)\) 的取值区域是顶点为 \((0,0)\)，\((1,0)\)，\((0,1)\) 的三角形. 代入顶点 \((1,0)\) 和 \((0,1)\)，题设恒成立的必要条件分别为 \(a\leqslant1\) 和 \(b\leqslant1\). 反之，若 \(0\leqslant a\leqslant1\)，\(0\leqslant b\leqslant1\)，则对任意满足题设条件的 \((x,y)\)，有 \(ax+by\leqslant x+y\leqslant1\)，所以这两个条件也充分. 因而点 \(P(a,b)\) 的区域为 \(0\leqslant a\leqslant1\)，\(0\leqslant b\leqslant1\) 围成的正方形，面积为 \(1\)，选择 C.
\end{solution}

\section{填空题}

\begin{problem}
已知函数 \( f(x) = x^2 + |x - 2| \)，则 \( f(1) = \)\fillinblank{}.
\end{problem}

\begin{answer}
\(2\)
\end{answer}

\begin{solution}
把 \(x=1\) 代入函数式，得 \(f(1)=1^2+|1-2|=1+1=2\)，所以填 \(2\).
\end{solution}

\begin{problem}
若 \(\sin \left(\frac{\pi}{2} + \theta\right) = \frac{3}{5}\)，则 \(\cos 2\theta =\)\fillinblank{}.
\end{problem}

\begin{answer}
\(-\frac{7}{25}\)
\end{answer}

\begin{solution}
由诱导公式 \(\sin\left(\frac\pi2+\theta\right)=\cos\theta\)，得 \(\cos\theta=\frac35\). 再由二倍角公式 \(\cos2\theta=2\cos^2\theta-1\)，得 \(\cos2\theta=2\left(\frac35\right)^2-1=\frac{18}{25}-\frac{25}{25}=-\frac7{25}\)，所以填 \(-\frac7{25}\).
\end{solution}

\begin{problem}
已知 \(F_{1}\)，\(F_{2}\) 为椭圆 \(\frac{x^{2}}{25} + \frac{y^{2}}{9} = 1\) 的两个焦点，过 \(F_{1}\) 的直线交椭圆于 \(A\)，\(B\) 两点. 若 \(|F_{2}A| + |F_{2}B| = 12\)，则 \(|AB| = \)\fillinblank{}.
\end{problem}

\begin{answer}
\(8\)
\end{answer}

\begin{solution}
由椭圆方程可知长半轴为 \(5\)，所以椭圆的焦点定义给出，对任意椭圆上的点 \(X\)，都有 \(|F_1X|+|F_2X|=10\). 因而 \(|F_2A|=10-|F_1A|\)，\(|F_2B|=10-|F_1B|\). 又因为直线 \(AB\) 经过椭圆内部的焦点 \(F_1\)，所以 \(F_1\) 在线段 \(AB\) 上，从而 \(|F_1A|+|F_1B|=|AB|\). 代入已知条件，得 \(12=|F_2A|+|F_2B|=20-|AB|\)，解得 \(|AB|=8\)，所以填 \(8\).
\end{solution}

\begin{problem}
在 \(\triangle ABC\) 中，角 \(A\)，\(B\)，\(C\) 所对的边分别为 \(a\)，\(b\)，\(c\)，若 \((\sqrt{3}b - c)\cos A = a\cos C\)，则 \(\cos A = \)\fillinblank{}.
\end{problem}

\begin{answer}
\(\frac{\sqrt{3}}{3}\)
\end{answer}

\begin{solution}
由正弦定理，存在 \(R>0\)，使 \(a=2R\sin A\)，\(b=2R\sin B\)，\(c=2R\sin C\). 将这些关系代入题设并约去正数 \(2R\)，得 \((\sqrt3\sin B-\sin C)\cos A=\sin A\cos C\). 移项后，\(\sqrt3\sin B\cos A=\sin C\cos A+\sin A\cos C=\sin(A+C)\). 因为三角形内角满足 \(A+B+C=\pi\)，所以 \(\sin(A+C)=\sin(\pi-B)=\sin B\). 又 \(0<B<\pi\)，故 \(\sin B>0\)，可以约去 \(\sin B\)，得到 \(\sqrt3\cos A=1\). 因此 \(\cos A=\frac1{\sqrt3}=\frac{\sqrt3}{3}\)，所以填 \(\frac{\sqrt3}{3}\).
\end{solution}

\begin{problem}
如图，已知球 \(O\) 的球面上四点 \(A\)，\(B\)，\(C\)，\(D\)，\(DA \perp\) 平面 \(ABC\)，\(AB \perp BC\)，\(DA = AB = BC = \sqrt{3}\)，则球 \(O\) 的体积等于\fillinblank{}.

\bitmapfigure[max width=0.35\linewidth]{img/2008/zhejiang_liberal/q15_fig1.png}
\end{problem}

\begin{answer}
\(\frac{9\pi}{2}\)
\end{answer}

\begin{solution}
由 \(AB\perp BC\) 及 \(AB=BC=\sqrt3\)，得 \(AC^2=AB^2+BC^2=6\). 又因为 \(DA\perp\) 平面 \(ABC\)，所以 \(DA\perp AC\)，从而 \(\triangle DAC\) 是直角三角形，且 \(DC^2=DA^2+AC^2=3+6=9\)，即 \(DC=3\). 另外，\(DA\perp BC\) 且 \(AB\perp BC\)，所以 \(BD\perp BC\)，从而 \(\triangle DBC\) 也是直角三角形. 因此 \(A\)，\(B\) 都在以 \(DC\) 为直径的球面上；又因 \(A\)，\(B\)，\(C\) 不共线且 \(D\notin\) 平面 \(ABC\)，所以 \(A\)，\(B\)，\(C\)，\(D\) 不共面，过这四点的球唯一，从而给定球面就是以 \(DC\) 为直径的球面. 球半径为 \(r=\frac{DC}{2}=\frac32\)，所以体积为 \(V=\frac43\pi r^3=\frac43\pi\left(\frac32\right)^3=\frac{9\pi}{2}\)，所以填 \(\frac{9\pi}{2}\).
\end{solution}

\begin{problem}
已知 \(\bs{a}\) 是平面内的单位向量，若向量 \(\bs{b}\) 满足 \(\bs{b} \cdot (\bs{a} - \bs{b}) = 0\)，则 \(|\bs{b}|\) 的取值范围是\fillinblank{}.
\end{problem}

\begin{answer}
\([0,1]\)
\end{answer}

\begin{solution}
由 \(\bs{b}\cdot(\bs{a}-\bs{b})=0\)，得 \(\bs{a}\cdot\bs{b}=|\bs{b}|^2\). 令 \(t=|\bs{b}|\geqslant0\)，因为 \(|\bs{a}|=1\)，由柯西不等式得 \(t^2=\bs{a}\cdot\bs{b}\leqslant|\bs{a}|\,|\bs{b}|=t\)，所以 \(0\leqslant t\leqslant1\). 反过来，任取 \(t\in[0,1]\)，在该平面内取与 \(\bs{a}\) 垂直的单位向量 \(\bs{c}\)，令 \(\bs{b}=t\left(t\bs{a}+\sqrt{1-t^2}\,\bs{c}\right)\). 由于 \(\bs{a}\perp\bs{c}\) 且 \(|\bs{a}|=|\bs{c}|=1\)，有 \(\left|t\bs{a}+\sqrt{1-t^2}\,\bs{c}\right|=1\)，从而 \(|\bs{b}|=t\). 同时 \(\bs{a}\cdot\bs{b}=t^2=|\bs{b}|^2\)，满足题设条件. 因此所有 \(t\in[0,1]\) 都能取得，所求范围为 \([0,1]\).
\end{solution}

\begin{problem}
用 \(1\)，\(2\)，\(3\)，\(4\)，\(5\)，\(6\) 组成六位数（没有重复数字），要求任何相邻两个数字的奇偶性不同，且 \(1\) 和 \(2\) 相邻，这样的六位数的个数是\fillinblank{}. （用数字作答）
\end{problem}

\begin{answer}
40
\end{answer}

\begin{solution}
六个数字中奇数有 \(1\)，\(3\)，\(5\) 三个，偶数有 \(2\)，\(4\)，\(6\) 三个. 相邻数字奇偶性不同，故奇偶位置只有“奇、偶、奇、偶、奇、偶”和“偶、奇、偶、奇、偶、奇”两种排列.

先考虑第一种位置模式. 数字 \(1\) 放在第 1，3，5 个位置时，分别有 \(1\)，\(2\)，\(2\) 个相邻的偶数位置可放数字 \(2\)，所以数字 \(1\)，\(2\) 相邻的放法有 \(1+2+2=5\) 种. 放好 \(1\)，\(2\) 后，剩下两个奇数和两个偶数各自排列，有 \(2!\times2!=4\) 种，故共有 \(5\times2!\times2!=20\) 个.

第二种位置模式中，数字 \(1\) 放在第 2，4，6 个位置时，可选的相邻偶数位置数为 \(2\)，\(2\)，\(1\)，同样有 \(5\times2!\times2!=20\) 个. 两种模式互斥，所以总数为 \(20+20=40\)，填 \(40\).
\end{solution}

\section{解答题}

\begin{problem}
已知数列 \(\{x_{n}\}\) 的首项 \(x_{1} = 3\)，通项 \(x_{n} = 2^{n}p + nq\)（\(n \in \N^{*}\)，\(p\)，\(q\) 为常数），且 \(x_{1}\)，\(x_{4}\)，\(x_{5}\) 成等差数列. 求：

\begin{enumerate}
\item \(p\)，\(q\) 的值；
\item 数列 \(\{x_{n}\}\) 前 \(n\) 项和 \(S_{n}\) 的公式.
\end{enumerate}
\end{problem}

\begin{solution}
\begin{enumerate}
\item 由通项公式，得 \(x_1=2p+q\)，\(x_4=16p+4q\)，\(x_5=32p+5q\). 由首项条件 \(x_1=3\)，得 \(2p+q=3\). 因为 \(x_1\)，\(x_4\)，\(x_5\) 成等差数列，所以中间项满足 \(x_1+x_5=2x_4\)，即 \((2p+q)+(32p+5q)=2(16p+4q)\). 化简得 \(p=q\). 与 \(2p+q=3\) 联立，得 \(3p=3\)，所以 \(p=q=1\).

\item 由 \(p=q=1\)，得对任意 \(k\in\N^*\)，\(x_k=2^k+k\). 因此
\[
S_n=\sum_{k=1}^{n}x_k=\sum_{k=1}^{n}2^k+\sum_{k=1}^{n}k=(2^{n+1}-2)+\frac{n(n+1)}{2}
\]
故 \(S_n=2^{n+1}-2+\dfrac{n(n+1)}{2}\).
\end{enumerate}
\end{solution}

\begin{problem}
一个袋中装有大小相同的黑球，白球和红球，已知袋中共有 \(10\) 个球，从中任意摸出 \(1\) 个球，得到黑球的概率是 \(\frac{2}{5}\)；从中任意摸出 \(2\) 个球，至少得到 \(1\) 个白球的概率是 \(\frac{7}{9}\). 求：

\begin{enumerate}
\item 从中任意摸出 \(2\) 个球，得到的 \(2\) 个球都是黑球的概率；
\item 袋中白球的个数.
\end{enumerate}
\end{problem}

\begin{solution}
设袋中黑球、白球、红球的个数分别为 \(b\)，\(w\)，\(r\). 由题意，\(b+w+r=10\)，且从中任取一个球得到黑球的概率为 \(\frac{b}{10}=\frac25\)，所以 \(b=4\)，从而 \(w+r=6\).

\begin{enumerate}
\item 任取两个球的所有结果等可能，共有 \(\mathrm C_{10}^{2}=45\) 种. 两个球都是黑球的结果有 \(\mathrm C_{4}^{2}=6\) 种，所以所求概率为 \(\frac{\mathrm C_4^2}{\mathrm C_{10}^2}=\frac6{45}=\frac2{15}\).

\item 设事件 \(E\) 为任取两个球至少有一个白球，则其对立事件 \(\overline E\) 为两个球中没有白球. 因而 \(P(\overline E)=1-P(E)=1-\frac79=\frac29\). 没有白球时，两个球必须从 \(10-w\) 个非白球中取出，所以 \(\frac{\mathrm C_{10-w}^{2}}{\mathrm C_{10}^{2}}=\frac29\)，即 \(\mathrm C_{10-w}^{2}=10\). 于是 \((10-w)(9-w)=20\)，化为 \(w^2-19w+70=0\)，解得 \(w=5\) 或 \(w=14\). 因为白球个数满足 \(0\leqslant w\leqslant10\)，所以 \(w=5\).
\end{enumerate}
所以袋中白球有 \(5\) 个.
\end{solution}

\begin{problem}
如图，矩形 \(ABCD\) 和梯形 \(BEFC\) 所在平面互相垂直，\(BE \parallel CF\)，\(\angle BCF = \angle CEF = 90^\circ\)，\(AD = \sqrt{3}\)，\(EF = 2\).

\begin{enumerate}
\item 求证：\(AE \parallel\) 平面 \(DCF\)；
\item 当 \(AB\) 的长为何值时，二面角 \(A - EF - C\) 的大小为 \(60^{\circ}\)？
\end{enumerate}

\bitmapfigure[max width=0.44\linewidth]{img/2008/zhejiang_liberal/q20_fig1.png}
\end{problem}

\begin{solution}
\begin{enumerate}
\item 在梯形 \(BEFC\) 所在平面内，过 \(E\) 作 \(EG\perp CF\)，垂足为 \(G\). 因为 \(BE\parallel CF\) 且 \(BC\perp CF\)，而 \(EG\) 也垂直于 \(CF\)，所以 \(EG\parallel BC\). 又因 \(G\in CF\)，有 \(CG\parallel BE\)，故四边形 \(BCGE\) 是平行四边形，从而 \(EG=BC\). 矩形 \(ABCD\) 给出 \(AD\parallel BC\) 且 \(AD=BC\)，所以 \(AD\parallel EG\) 且 \(AD=EG\)，四边形 \(ADGE\) 是平行四边形，进而 \(AE\parallel DG\). 由于 \(D\)，\(G\)，\(C\)，\(F\) 共面且 \(G\in CF\)，所以 \(DG\subset\) 平面 \(DCF\). 又平面 \(ABCD\) 与平面 \(DCF\) 的交线为 \(DC\)，而 \(A\notin DC\)，故 \(A\notin\) 平面 \(DCF\). 由直线与平面平行的判定，\(AE\parallel\) 平面 \(DCF\).

\item 取梯形 \(BEFC\) 所在平面为 \(z=0\)，以 \(B\) 为原点、\(BC\) 所在直线为 \(x\) 轴，建立空间直角坐标系. 由 \(AD=BC=\sqrt3\)，以及 \(BE\parallel CF\)，\(BC\perp CF\)，可设 \(B=(0,0,0)\)，\(C=(\sqrt3,0,0)\)，\(E=(0,m,0)\)，\(F=(\sqrt3,n,0)\)，其中 \(m=BE>0\)，\(n=CF>0\). 矩形所在平面与梯形所在平面垂直，且矩形中 \(AB\perp BC\)，所以 \(AB\perp\) 梯形所在平面. 设 \(h=AB>0\)，则可取 \(A=(0,0,h)\)，\(D=(\sqrt3,0,h)\).

由 \(EF=2\)，得 \(3+(n-m)^2=4\)，即 \((n-m)^2=1\). 又因为 \(\angle CEF=90^\circ\)，所以
\[
\overrightarrow{EC}\cdot\overrightarrow{EF}=(\sqrt3,-m,0)\cdot(\sqrt3,n-m,0)=3-m(n-m)=0
\]
因此 \(m(n-m)=3>0\). 结合 \(m>0\) 和 \((n-m)^2=1\)，得 \(n-m=1\)，从而 \(m=3\)，\(n=4\).

平面 \(CEF\) 的一个法向量为 \(\bs{n}_1=(0,0,1)\). 平面 \(AEF\) 的两个方向向量可取 \(\overrightarrow{AE}=(0,3,-h)\) 和 \(\overrightarrow{AF}=(\sqrt3,4,-h)\)，所以其法向量可取
\[
\bs{n}_2=\overrightarrow{AE}\times\overrightarrow{AF}=(0,3,-h)\times(\sqrt3,4,-h)=(h,-\sqrt3h,-3\sqrt3)
\]
设二面角 \(A-EF-C\) 的大小为 \(\theta\). 二面角的平面角等于两侧平面法向量的锐角夹角，因此
\[
\cos\theta=\frac{|\bs{n}_1\cdot\bs{n}_2|}{|\bs{n}_1|\,|\bs{n}_2|}=\frac{3\sqrt3}{\sqrt{4h^2+27}}
\]
当 \(\theta=60^\circ\) 时，
\(\frac{3\sqrt3}{\sqrt{4h^2+27}}=\frac12\)，\(4h^2+27=108\)，\(h^2=\frac{81}{4}\).
因为 \(h=AB>0\)，所以 \(AB=\frac92\). 因而当 \(AB=\frac92\) 时，二面角 \(A-EF-C\) 的大小为 \(60^\circ\).
\end{enumerate}
\end{solution}

\begin{problem}
已知 \(a\) 是实数，函数 \(f(x) = x^2 (x - a)\).

\begin{enumerate}
\item 若 \(f'(1)=3\)，求 \(a\) 的值及曲线 \(y=f(x)\) 在点 \((1,f(1))\) 处的切线方程；
\item 求 \(f(x)\) 在区间 \([0,2]\) 上的最大值.
\end{enumerate}
\end{problem}

\begin{solution}
\begin{enumerate}
\item 由 \(f(x)=x^{3}-ax^{2}\)，得 \(f'(x)=3x^{2}-2ax\). 由 \(f'(1)=3\)，得 \(3-2a=3\)，所以 \(a=0\). 此时 \(f(x)=x^{3}\)，点 \((1,f(1))=(1,1)\)，且切线斜率为 \(f'(1)=3\). 因此切线方程为 \(y-1=3(x-1)\)，即 \(y=3x-2\).

\item 第二问中的 \(a\) 仍为任意实数，不能把第一问附加的条件 \(f'(1)=3\) 当作本问条件. 对 \(f(x)=x^3-ax^2\) 求导，得 \(f'(x)=x(3x-2a)\). 在区间 \([0,2]\) 内，临界点除端点 \(x=0\) 外为 \(x=\frac{2a}{3}\). 当 \(a\leqslant0\) 时，\(f'(x)\geqslant0\)（\(0\leqslant x\leqslant2\)），函数单调递增，最大值为 \(f(2)=8-4a\). 当 \(a\geqslant3\) 时，\(f'(x)\leqslant0\)（\(0\leqslant x\leqslant2\)），函数单调递减，最大值为 \(f(0)=0\). 当 \(0<a<3\) 时，函数在 \([0,\frac{2a}{3}]\) 上递减，在 \([\frac{2a}{3},2]\) 上递增，因此最大值只能在两个端点取得. 比较 \(f(0)=0\) 与 \(f(2)=8-4a\)：当 \(0<a\leqslant2\) 时最大值为 \(8-4a\)，当 \(2<a<3\) 时最大值为 \(0\). 综合各情形，
\[
\max_{x\in[0,2]}f(x)=
\begin{cases}
8-4a,&a\leqslant2,\\
0,&a>2.
\end{cases}
\]
\end{enumerate}
\end{solution}

\begin{problem}
已知曲线 \(C\) 是到点 \(P\left(-\frac{1}{2}, \frac{3}{8}\right)\) 和到直线 \(y = -\frac{5}{8}\) 距离相等的点的轨迹. \(l\) 是过点 \(Q(-1, 0)\) 的直线，\(M\) 是 \(C\) 上（不在 \(l\) 上）的动点，\(A\)，\(B\) 在 \(l\) 上，\(MA \perp l\)，\(MB \perp x\) 轴（如图）.

\begin{enumerate}
\item 求曲线 \(C\) 的方程；
\item 求出直线 \(l\) 的方程，使得 \(\frac{|QB|^{2}}{|QA|}\) 为常数.
\end{enumerate}

\bitmapfigure[max width=0.46\linewidth]{img/2008/zhejiang_liberal/q22_fig1.png}
\end{problem}

\begin{solution}
\begin{enumerate}
\item 设曲线 \(C\) 上任意一点为 \(M(x,y)\). 点 \(M\) 到点 \(P\) 的距离为 \(\sqrt{\left(x+\frac12\right)^2+\left(y-\frac38\right)^2}\)，到直线 \(y=-\frac58\) 的距离为 \(\left|y+\frac58\right|\). 由题意，
\[
\left(x+\frac12\right)^2+\left(y-\frac38\right)^2=\left(y+\frac58\right)^2
\]
展开得 \(x^2+x+\frac14+y^2-\frac34y+\frac9{64}=y^2+\frac54y+\frac{25}{64}\)，整理为 \(x^2+x=2y\). 所以曲线 \(C\) 的方程为 \(y=\frac{x^2+x}{2}\).

\item 若 \(l\perp x\) 轴，则 \(l\) 为竖直直线 \(x=-1\). 对于 \(C\) 上不在 \(l\) 上的点 \(M\)，过 \(M\) 的竖直线与 \(l\) 平行，不能交于点 \(B\)，故 \(l\) 不是竖直线. 设其斜率为 \(k\)，则 \(l:y=k(x+1)\). 由 \(MB\perp x\) 轴，得 \(B=(x,k(x+1))\)，因此 \(|QB|^2=(1+k^2)(x+1)^2\).

取 \((1,k)\) 为 \(l\) 的方向向量. 因为 \(MA\perp l\)，点 \(A\) 是 \(M\) 在直线 \(l\) 上的正投影，故 \(\overrightarrow{QA}\) 是 \(\overrightarrow{QM}=(x+1,y)\) 在 \((1,k)\) 上的投影：
\(\overrightarrow{QA}=\frac{x+1+ky}{1+k^2}(1,k)\)，\(|QA|=\frac{|x+1+ky|}{\sqrt{1+k^2}}\).
代入曲线方程，得 \(x+1+ky=(x+1)\left(1+\frac{kx}{2}\right)\)，从而
\[
\frac{|QB|^2}{|QA|}=(1+k^2)^{3/2}\frac{|x+1|}{\left|1+\frac{kx}{2}\right|}
\]
若 \(k=0\)，上式为 \(|x+1|\)，随曲线上的动点变化，不可能为常数. 因此 \(k\ne0\). 对于 \(k\ne0\)，若比值对所有允许的动点恒定，则在除去有限个交点后，对无穷多个 \(x\) 有 \(\frac{|x+1|}{|1+\frac{kx}{2}|}\) 为常数. 两边平方并交叉相乘，得到的两个二次多项式在无穷多个点相等，因而恒等；所以两个一次式 \(x+1\) 与 \(1+\frac{kx}{2}\) 必须成比例. 比较它们的零点，得 \(-1=-\frac2k\)，所以 \(k=2\).

还需排除分母为零而 \(M\notin l\) 的情形. 当 \(k\ne2\) 时，取 \(x=-\frac2k\)，曲线 \(C\) 上相应点的纵坐标为 \(\frac{2-k}{k^2}\)，而直线 \(l\) 在该横坐标处的纵坐标为 \(k-2\). 若该点在 \(l\) 上，则
\[
\frac{2-k}{k^2}=k-2
\]
即 \((k-2)(k^2+1)=0\)，这与 \(k\ne2\) 矛盾. 所以该点不在 \(l\) 上，却使 \(|QA|=0\)，题设比值无意义，故 \(k\ne2\) 的情形均不合适. 当 \(k=2\) 时，分母为零只发生在 \(x=-1\)，对应点 \(Q\)，而题设已排除 \(M\in l\). 此时 \(1+\frac{kx}{2}=x+1\)，所以 \(\frac{|QB|^2}{|QA|}=5^{3/2}=5\sqrt5\) 为常数. 因而所求直线为 \(l:y=2x+2\).
\end{enumerate}
\end{solution}
